% !TeX root = ../main.tex
% System Model and Problem Formulation (input into \section in main.tex).

The main symbols are summarized in Table~\ref{tab:notation}.

\begin{table}[t]
  \centering
  \caption{Summary of notation.}
  \label{tab:notation}
  \begin{tabular}{ll}
    \toprule
    Symbol & Meaning \\
    \midrule
    $K,\;\mathcal{K}$ & number and set of ground sensors \\
    $M$ & swarm size (number of UAVs) \\
    $\mathcal{S},\;S_{\max}$ & candidate sites, station budget \\
    $x_s\in\{0,1\}$ & open a station at candidate site $s$ \\
    $c_s,\;C_{\mathrm{tot}}$ & ports at station $s$, total port budget \\
    $\mathbf{q}_m(t)$ & trajectory of UAV $m$ \\
    $a_m$ & station assigned to UAV $m$ \\
    $A_k(t),\;\bar A$ & AoI of sensor $k$, network peak AoI (objective) \\
    $V,\;H$ & cruise speed, altitude \\
    $E_{\max},\;\eta$ & battery, reserve fraction \\
    $\tau_{\mathrm{fly}},\;\tau_{\mathrm{charge}}$ & flight budget, charge time \\
    $T_{\mathrm{up}}$ & operating time between charges (incl. travel) \\
    $N_s$ & UAVs assigned to station $s$ \\
    $\lambda,\;\mu$ & charge-request rate, service rate \\
    $\Wq,\;t_{\mathrm{loop}}$ & mean queue wait, sub-field revisit period \\
    \bottomrule
  \end{tabular}
\end{table}

\subsection{Network model}
\label{subsec:network}
We consider a swarm data-collection system, illustrated in Fig.~\ref{fig:system}.
A set $\mathcal{K}=\{1,\dots,K\}$ of ground sensors is deployed over a square
field, with sensor $k$ located at $\mathbf{w}_k\in\mathbb{R}^2$. A swarm of
$M$ rotary-wing UAVs, indexed by $\mathcal{M}=\{1,\dots,M\}$, flies at a fixed
altitude $H$ with horizontal position $\mathbf{q}_m(t)\in\mathbb{R}^2$ over
discrete slots $t=1,\dots,T$ of length $\delta_t$. Each UAV holds a battery with
state of charge $e_m(t)\in[0,E_{\max}]$. A set
$\mathcal{S}=\{1,\dots,S\}$ of candidate charging-station sites is given, with
site $s$ at $\mathbf{b}_s\in\mathbb{R}^2$; the system may build a station and
provision it with charging ports.

Freshness is measured by the AoI~\cite{kaul2012aoi}. Let
$u_k(t)$ be the generation time of the freshest update from sensor $k$ that has
been delivered by time $t$. The instantaneous AoI is
\begin{equation}
  \Aoi_k(t) = t - u_k(t),
  \label{eq:aoi}
\end{equation}
which grows at unit rate and drops on each successful collection: when a UAV
collects at time $\tau$ a packet generated at $g$, $u_k$ jumps to $g$ and
$\Aoi_k(\tau^{+})=\tau-g$. Following the base setting, the network objective is
the peak AoI across sensors,
% ⛔ DOI o vong sua R1, theo diem 1 cua phan bien. Ban truoc dinh nghia \bar{A} la
% worst-case theo QUY DAO MAU (limsup) nhung moi bieu thuc sau do lai dung Wq TRUNG
% BINH. Hai dai luong khac nhau mang chung mot ten. Da DO xem thay the co duoc khong:
% ti so cho o phan vi 99 tren cho trung binh len toi 82,9 lan o che do cap du dat,
% nen KHONG duoc. Nay dinh nghia la KY VONG, tuc dung thu that su duoc toi uu.
% ⚠ Cot mot cot cua IEEEtran chi rong ~252pt: viet ca cong thuc va cum giai thich
% tren MOT dong thi tran 67pt. Tach cum giai thich xuong van xuoi.
let $t_{k,1}<t_{k,2}<\cdots$ be the epochs at which an update from sensor $k$ is
delivered, so that $\Aoi_k(t_{k,i}^-)$ is the age reached just before the $i$-th
delivery, i.e. the \emph{peak} of cycle $i$. The objective is the expected peak age of
the worst sensor,
\begin{equation}
  \bar{A} \;=\; \max_{k\in\mathcal{K}}\ \limsup_{N\to\infty}\
  \frac{1}{N}\sum_{i=1}^{N}\Aoi_k\!\left(t_{k,i}^{-}\right),
  \label{eq:peakaoi}
\end{equation}
which we seek to minimize.

Three points about this definition, since an earlier version of this work was
inconsistent about it. First, \eqref{eq:peakaoi} is an \emph{expectation over cycles},
not a sample-path worst case: the earlier form $\max_k\limsup_t \Aoi_k(t)$ is withdrawn,
and \eqref{eq:P0} below optimizes \eqref{eq:peakaoi} rather than restating a different
functional. Second, \eqref{eq:peakaoi} is the average of the \emph{peaks}, not the
time-average of the age; those differ by a factor close to two, and we measure rather
than assert it. Across the $\rPaoiConfigs$ configurations of
Section~\ref{sec:metric-gap} the ratio of the expected peak to the time-average age
is $\rPaoiRatioLo$ to $\rPaoiRatioHi$, because the age sawtooths from zero to the
revisit interval and the cycle length varies little (coefficient of variation at most
$\rPaoiCvHi$). Reporting a time-average under the name ``peak AoI'' would therefore
halve every number in this paper. Third, using the \emph{mean} charging wait $\Wq$
inside \eqref{eq:peakaoi} is legitimate precisely because \eqref{eq:peakaoi} is itself
an expectation; it would not be legitimate under the withdrawn sample-path form, and
Section~\ref{sec:metric-gap} quantifies that gap: in the well-provisioned regime this
paper recommends, the largest observed wait exceeds the mean by up to $\rMaxHi\times$. A UAV collects from sensor $k$ while
$\|\mathbf{q}_m(t)-\mathbf{w}_k\|\le r_c$.

\subsection{Rotary-wing energy}
\label{subsec:energy}
The charging cadence is not assumed constant: it emerges from how fast a UAV
drains its battery in flight. For a rotary-wing UAV cruising at horizontal speed
$V$, the propulsion power follows the closed-form model of Zeng~\emph{et
al.}~\cite{zeng2019energy},
\begin{equation}
\begin{split}
  P(V) = {}& P_0\!\left(1+\frac{3V^2}{U_{\mathrm{tip}}^2}\right)
       + P_i\!\left(\sqrt{1+\frac{V^4}{4v_0^4}}-\frac{V^2}{2v_0^2}\right)^{\!1/2}\\
       &{}+ \tfrac{1}{2}d_0\rho s A\,V^3,
\end{split}
  \label{eq:power}
\end{equation}
where $P_0$ and $P_i$ are the blade-profile and induced power constants,
$U_{\mathrm{tip}}$ the rotor tip speed, $v_0$ the mean rotor induced velocity in
hover, and $d_0,\rho,s,A$ the fuselage drag ratio, air density, rotor solidity,
and disc area. Hovering during collection consumes $P(0)=P_0+P_i$. The charge
evolves as $e_m(t+1)=e_m(t)-P(V_m(t))\,\delta_t$ while flying or collecting, and
increases while docked at a charging port.

Only a usable fraction of the pack is spent between charges, keeping a reserve
$\eta$, so $E_{\mathrm{use}}=(1-\eta)E_{\max}$. Given a patrol loop of flight
time $t_{\mathrm{fly}}^{\mathrm{loop}}$ and hover-collection time
$t_{\mathrm{col}}^{\mathrm{loop}}$, the energy drawn per loop is
$e_{\mathrm{loop}} = P(V)\,t_{\mathrm{fly}}^{\mathrm{loop}}
+ P(0)\,t_{\mathrm{col}}^{\mathrm{loop}}$, so a UAV can fly
$n_{\mathrm{loop}}=\lfloor E_{\mathrm{use}}/e_{\mathrm{loop}}\rfloor$ loops before
recharging, giving a physical flight budget between charges
\begin{equation}
  \tau_{\mathrm{fly}} = n_{\mathrm{loop}}\, t_{\mathrm{loop}},
  \qquad
  t_{\mathrm{loop}} = t_{\mathrm{fly}}^{\mathrm{loop}}+t_{\mathrm{col}}^{\mathrm{loop}}.
  \label{eq:taufly}
\end{equation}
A recharge-feasibility condition requires each UAV to retain enough charge to
reach its assigned open station at all times:
$e_m(t)\ge E_{\mathrm{reach}}(\mathbf{q}_m(t),\mathbf{b}_{a_m})$.

\subsection{Air-to-ground collection}
\label{subsec:comms}
Collection dwell is derived from a physical link rather than assumed
instantaneous. With a line-of-sight, free-space channel of reference gain
$\beta_0$ at $1$~m, a sensor at horizontal distance $d$ is seen at slant distance
$\sqrt{H^2+d^2}$, and the achievable rate is
\begin{equation}
  R(d) = B\,\log_2\!\left(1+\frac{p_{\mathrm{tx}}\,\beta_0}
       {(H^2+d^2)\,N_0 B}\right),
  \label{eq:rate}
\end{equation}
with bandwidth $B$, sensor transmit power $p_{\mathrm{tx}}$, and noise power
spectral density $N_0$. Collecting a status blob of $L$ bits at the
closest-approach distance therefore takes a dwell time
$t_{\mathrm{col}}(d)=L/R(d)$, which enters $t_{\mathrm{loop}}$ in
\eqref{eq:taufly} and hence the revisit period.

\subsection{Finite-source charging queue}
\label{subsec:queue}
The load-bearing coupling of the swarm is contention for charging ports. A
station $s$ serving the $N_s$ UAVs assigned to it with $c_s$ parallel ports is
\emph{not} an open queue: only $N_s$ UAVs exist, so the request rate falls as
more of them are already at the station. The correct model is the finite-source
(machine-repair) queue $M/M/c//N_s$~\cite{grossharris}.

Each assigned UAV, while patrolling, requests a charge after a mean operating
time, at rate $\lambda=1/T_{\mathrm{up}}$; each port serves at
$\mu=1/\tau_{\mathrm{charge}}$. The operating time between two charge requests is
the whole time a UAV spends away from the station in a cycle, so it must include
the round-trip travel to the station and not the flight budget alone:
\begin{equation}
  T_{\mathrm{up}} = \tau_{\mathrm{fly}} + \frac{2\,d(m,s)}{V},
  \label{eq:uptime}
\end{equation}
where $d(m,s)=\|\mathbf{q}_m-\mathbf{b}_s\|$. Omitting the travel term overstates $\lambda$ and overpredicts the wait by roughly a
factor of two.

\emph{Delivery semantics.} Two reviewers asked us to state when a collected update
counts as \emph{delivered}, since the age is defined by delivery while the system model
describes collection. We assume an ideal backhaul: the instant a UAV collects an update
it is available at the destination, so collection time equals delivery time and no
backhaul or data-carrying delay enters the age. This is an idealization, appropriate
only when a persistent air-to-ground or satellite backhaul is present. A
deliver-on-return model, in which the UAV carries the data to a ground node, would add
the return flight to every peak and would change \eqref{eq:peak_m}; the placement and
capacity structure we study would survive, but the absolute ages would not. We state
the assumption here rather than leave it implicit.

Let $n\in\{0,\dots,N_s\}$ be the number of UAVs at the station. The birth-death
chain has state-dependent rates
\begin{equation}
  n\!\to\! n{+}1:\ (N_s-n)\lambda,
  \qquad
  n\!\to\! n{-}1:\ \min(n,c_s)\mu,
  \label{eq:rates}
\end{equation}
where the up-rate $(N_s-n)\lambda$ vanishes as $n\to N_s$: this closed-population
term is exactly why the queue cannot blow up. The stationary distribution is
\begin{equation}
  p_n = p_0\prod_{k=0}^{n-1}\frac{(N_s-k)\,\lambda}{\min(k{+}1,c_s)\,\mu},
  \qquad
  \sum_{n=0}^{N_s} p_n = 1.
  \label{eq:pn}
\end{equation}
The mean number waiting, the effective throughput, and, by Little's law, the mean
queue wait are
\begin{align}
  L_q &= \sum_{n=0}^{N_s}\max(0,\,n-c_s)\,p_n,
  \label{eq:lqlam}\\
  \lambda_{\mathrm{eff}} &= \sum_{n=0}^{N_s}(N_s-n)\,\lambda\,p_n,
  \nonumber\\
  \Wq(N_s,c_s) &= \frac{L_q}{\lambda_{\mathrm{eff}}}.
  \label{eq:wq}
\end{align}
Because the population is finite, no stability constraint (of the form
$\rho_s<1$) is needed: the wait \emph{saturates} rather than diverging. This
finite-source model replaces the open $M/M/c$ (Erlang-C) form, which treats
arrivals as an inexhaustible stream and, for the cycling swarm, overpredicts the
wait by roughly a factor of six. A discrete-event validation of the true cycling
UAVs, reported in Section~\ref{sec:exp}, shows that \eqref{eq:wq} tracks the
simulation to within a few percent and, on every configuration we tested, overestimates the real,
near-deterministic wait.

\subsection{Peak-AoI expression}
\label{subsec:aoi}
\label{subsec:peakaoi}
For a UAV $m$ that patrols a sub-field and leaves it to recharge, the expected revisit
interval seen by its sensors is one patrol period plus one charging excursion, and the
age of a sensor peaks at exactly that interval just before the next delivery:
\begin{equation}
  \Aoi_m = t_{\mathrm{loop},m}
         + \frac{2\,d(m,s)}{V}
         + \Wq(N_s,c_s)
         + \tau_{\mathrm{charge}}.
  \label{eq:peak_m}
\end{equation}
The first term is the patrol revisit period, the second the round trip to the
assigned station, the third the finite-source queue wait from \eqref{eq:wq} (the
term that couples the swarm), and the last the charging time. Because every term is a
mean, \eqref{eq:peak_m} is the expected cycle peak of \eqref{eq:peakaoi} for the
sensors served by UAV $m$; the network objective \eqref{eq:peakaoi} is the maximum of
\eqref{eq:peak_m} over the swarm. This is the quantity Algorithm~1 minimizes, and it
is the same functional that \eqref{eq:P0} states. The
coupling in $\Wq$ is the mechanism behind the non-monotone behavior of
$\bar{A}$ in the swarm size $M$: more UAVs shrink the coverage term but lengthen
the wait, up to the ceiling set by the finite population.

\subsection{Joint optimization problem}
\label{subsec:problem}
We jointly optimize the placement decisions $x_s\in\{0,1\}$ (build a station at
$s$), the capacities $c_s\in\mathbb{Z}_{\ge0}$ (ports per station), the UAV
trajectories $\{\mathbf{q}_m(\cdot)\}$, and the station assignment
$\{a_m\}$. The problem is
\begin{subequations}\label{eq:P0}
\begin{align}
  \text{(P0)}\quad
  \argmin_{\substack{\mathbf{x},\,\mathbf{c},\\ \{\mathbf{q}_m\},\,\{a_m\}}}\ \
    & \bar{A}\ \text{ as defined in }\eqref{eq:peakaoi}
    \label{eq:P0obj}\\
  \text{s.t.}\quad
    & e_m(t)\in[0,E_{\max}],\ \forall m,t,
    \label{eq:P0e1}\\
    & e_m(t)\ge E_{\mathrm{reach}}(\mathbf{q}_m(t),\mathbf{b}_{a_m}),\ \forall m,t,
    \label{eq:P0e2}\\
    & \Aoi_k(t)\ \text{evolves as in \eqref{eq:aoi}},\ \forall k,t,
    \label{eq:P0aoi}\\
    & \textstyle\sum\nolimits_{s} x_s \le S_{\max},\quad
      \sum\nolimits_{s} c_s \le C_{\mathrm{tot}},
    \label{eq:P0bud}\\
    & c_s \le C_{\mathrm{port}}\,x_s,\quad
      x_s\in\{0,1\},\ c_s\in\mathbb{Z}_{\ge0},
    \label{eq:P0cap}\\
    & \|\mathbf{q}_m(t{+}1)-\mathbf{q}_m(t)\|\le V_{\max}\delta_t,\ \forall m,t,
    \label{eq:P0kin}\\
    & a_m\in\{s:x_s=1\},\ \forall m,
    \label{eq:P0asg}\\
    & \|\mathbf{q}_m(t)-\mathbf{q}_{m'}(t)\|\ge d_{\min},\ \forall m\ne m',t.
    \label{eq:P0sep}
\end{align}
\end{subequations}
Constraint \eqref{eq:P0e2} is recharge feasibility, \eqref{eq:P0aoi} ties peak
AoI to collection events through \eqref{eq:peak_m} and the queue wait
\eqref{eq:wq}, \eqref{eq:P0bud} and \eqref{eq:P0cap} enforce the station and port
budgets, \eqref{eq:P0kin} is UAV kinematics, \eqref{eq:P0asg} restricts
assignment to open stations, and \eqref{eq:P0sep} is inter-UAV separation.

Problem (P0) mixes binary placement $\mathbf{x}$, integer capacity $\mathbf{c}$,
continuous trajectories, and a combinatorial assignment, and its objective embeds
the nonconvex finite-source wait \eqref{eq:wq} through the coupling
\eqref{eq:peak_m}. It is therefore a mixed-integer nonlinear program (MINLP) and
NP-hard, which motivates the block-coordinate decomposition (trajectory, queue
evaluation, capacity allocation, assignment, and outer placement) detailed in
Section~\ref{sec:solver}.
